\providecommand{\ZZ}{\mathbb{Z}}
\providecommand{\CC}{\mathbb{C}}
\providecommand{\PP}{\mathbb{P}}
\providecommand{\HH}{\mathbb{H}}
\providecommand{\kBKM}{\kappa_{\mathrm{BKM}}}
\providecommand{\kch}{\kappa_{\mathrm{ch}}}
\providecommand{\kcat}{\kappa_{\mathrm{cat}}}
\providecommand{\kfib}{\kappa_{\mathrm{fiber}}}
\providecommand{\Aut}{\mathrm{Aut}}
\providecommand{\Hilb}{\mathrm{Hilb}}
\providecommand{\Kum}{\mathrm{Kum}}
\providecommand{\Fix}{\mathrm{Fix}}
\providecommand{\Sp}{\mathrm{Sp}}
\providecommand{\ClaimStatusConjectured}{\textsc{[Conjectured]}}
\providecommand{\ClaimStatusProvedElsewhere}{\textsc{[ProvedElsewhere]}}

\section*{Wave~16 V1 --- non-CHL hosts for the $N\in\{5,7,8\}$ paramodular forms}

Wave~15 R4 sharpened the obstruction: for $N\in\{5,7,8\}$, the
CHL orbifold $(S\times E)/(\ZZ/N)$ collapses because $\varphi(N)\nmid 2$
rules out $\ZZ[\zeta_N]\hookrightarrow\mathrm{End}(E)$, the $N=7$ weight
$1/4$ requires a spin-double-cover character absent from the
$\mathcal{N}=(4,4)$ monodromy, and $\Delta_{0}^{(8)}$ is a weight-$0$
modular function with empty Weyl chamber. The Gritsenko--Cl\'ery
forms nevertheless exist; the question is what CY$_3$ hosts them.
The paramodular side of the universal Borcherds weight identity
$\kBKM(\Phi_N)=c_N(0)/2$ is indifferent to CHL --- it is a theorem of
singular-theta transfer. What breaks is the geometric realisation
$Z^X\leftrightarrow\Phi_N$. We propose the correct geometric side.

\subsection*{Cycle~1 --- Mongardi--Tari--Wandel Kummer-type CY$_3$ for $N=8$}

Mongardi--Tari--Wandel (arXiv:1610.06216, \emph{Symplectic involutions
of K3$^{[n]}$ type and Kummer$^{n}$ type manifolds}) classify
symplectic automorphisms of irreducible holomorphic-symplectic
manifolds of Kummer-$n$ type $\Kum^{n}(A)$, the generalised Kummer
fourfolds of Beauville--Bogomolov. Theorem~1.3 of op.~cit.\ exhibits
a symplectic order-$8$ automorphism on $\Kum^{3}(A)$ whose fixed
locus is the union of $8$ isolated points, lifting a translation-
and involution-compatible $\ZZ/8$ on the abelian surface $A$. The
associated CY$_3$ host is the \emph{quotient threefold}
\[
  X_{8}^{\mathrm{MTW}} \;:=\; \bigl(\Kum^{3}(A)\bigr) / (\ZZ/8),
\]
resolved by a small crepant resolution at the $8$ fixed points
(equivalently, by toric resolution of the local $\CC^{3}/(\ZZ/8)$
McKay quotient, giving $\kfib(X_{8}^{\mathrm{MTW}})=0$ by the Bridgeland--
King--Reid derived McKay correspondence). Because the Kummer locus
carries its own $\CC^{\times}$-action through $A$ translations, the
$N=8$ automorphism commutes with a free elliptic $\ZZ/2$, recovering
the Oberdieck--Pixton $N M\leq 8$ envelope \emph{without} requiring
$\ZZ[\zeta_{8}]$-CM on an elliptic curve.

\subsection*{Cycle~2 --- $N=5$ via Borcea--Voisin mirror of Reid's 95 families}

For $N=5$ no Kummer-type HK $6$-fold carries a symplectic order-$5$
action (Mongardi--Tari--Wandel Remark~2.4: Kummer lattices are
stable only under orders dividing $8$). The correct host is
instead the Borcea--Voisin CY$_3$
\[
  X_{5}^{\mathrm{BV}} \;:=\; \bigl(S_{5}\times E_{5}\bigr)/(\iota_{S}\times \iota_{E})
\]
where $S_{5}$ is a Nikulin involution-quotient K3 whose
transcendental lattice is invariant under a symplectic order-$5$
element $g_{5}\in\Aut_{s}(S_{5})$ with Nikulin fixed count
$\#\,\Fix(g_{5})=4$, and $\iota_{S}\times\iota_{E}$ is the
diagonal anti-symplectic involution constructing the Borcea--Voisin
threefold (Borcea \emph{K3 surfaces with involution and mirror pairs
of CY manifolds}, 1997; Voisin, \emph{Miroirs et involutions sur les
surfaces K3}, 1993). The elliptic factor $E_{5}$ carries an order-$2$
involution (not order~$5$): the discrepancy from CHL is that the
$\ZZ/5$ acts only on the K3 factor, and the transcendental information
$c_{5}(0)$ is read from the $g_{5}$-twisted elliptic genus of $S_{5}$
rather than the joint diagonal genus.

\subsection*{Cycle~3 --- twisted DT zeta $Z^{X}$ on the proposed hosts}

Oberdieck--Pixton arXiv:1808.03524 Conjecture~A computes
$Z^{S\times E}_{L}=-C/\Delta_{5}^{2}$ on the CHL slice. For
non-CHL hosts the analogous conjecture takes the form, with
lattice data $L\subset H^{2}(X;\ZZ)$ polarising the base and
$h_M=1$ (pure order-$N$ twist):
\begin{align*}
  Z^{X_{8}^{\mathrm{MTW}}}_{L,1}(p,q,y)
  \;&=\; -C_{8}\cdot \Phi_{0}^{(8)}(p,q,y)^{-2}\quad\text{at }k_{8}=0,\\
  Z^{X_{5}^{\mathrm{BV}}}_{L,1}(p,q,y)
  \;&=\; -C_{5}\cdot \bigl(\Delta_{1/2}^{(5)}(p,q,y)\bigr)^{-4}\quad\text{at }k_{5}=1/2,\\
  Z^{X_{7}^{\mathrm{?}}}_{L,1}(p,q,y)
  \;&=\; -C_{7}\cdot \bigl(\Delta_{1/4}^{(7)}(p,q,y)\bigr)^{-8}\quad\text{at }k_{7}=1/4,
\end{align*}
the exponents $-2, -4, -8$ chosen so that $2k_{N}\cdot|\text{exp}|=|c_{N}(0)|\cdot 2$
restores integer BKM multiplicities. The $N=7$ host is \emph{not}
yet identified in the Mongardi--Tari--Wandel census; Nikulin's
fixed count $\#\,\Fix(g_{7})=3$ combined with $\chi(\mathrm{K3}^{g_{7}})/3=0$
forces any $\ZZ/7$-orbifold threefold to be a singular limit of a
smooth CY$_3$ rather than a free quotient. This is the
open-geometry end of Wave~16 V1.

\subsection*{Cycle~4 --- $(Z^{X})^{-1/2}$ versus $\Delta_{w_N}^{(N)}$}

The universal Borcherds weight theorem (Borcherds 1998 Thm~13.3)
computes the weight of $\Phi_{N}$ from the EZ constant of its input
Jacobi form: $k_N=c_N(0)/2$ at $\phi_{0,1}^{\bullet}$-weight $0$
truncation. For the proposed hosts:
\begin{itemize}[leftmargin=1.8em,itemsep=0.3ex]
\item $N=8$: $\Delta_{0}^{(8)}$ has $c_8(0)=0$, so $\kBKM(\Phi_8)=0$.
  The GKM superalgebra $\mathfrak{g}_{\Phi_{8}}$ is abelian (empty
  real-root system); the $(Z^{X_{8}^{\mathrm{MTW}}})^{-1/2}=\Phi_{8}$
  match is trivially consistent at the weight-$0$ endpoint.
\item $N=5$: $\Delta_{1/2}^{(5)}$ has $c_5(0)=1$, so $\kBKM=1/2$.
  The BKM lift carries a half-integer Weyl character, realising a
  \emph{metaplectic} BKM superalgebra $\widetilde{\mathfrak{g}}_{\Phi_{5}}$.
  The identification $(Z^{X_{5}^{\mathrm{BV}}})^{-1/4}=\Phi_{5}$
  --- note the \emph{fourth} root --- is forced by the weight-$1/2$
  squaring: $\Phi_5^2$ is a weight-$1$ paramodular form with
  integer BKM data, and the fourth root of $Z^{X}$ produces the
  metaplectic square root.
\item $N=7$: $\Delta_{1/4}^{(7)}$ with $c_7(0)=1/2$ gives $\kBKM=1/4$,
  requiring an \emph{eighth-root} of $Z^{X}$; the spin-double-cover
  character $\chi^{\otimes 2}$ is geometrically the Cheeger--Simons
  differential character of an order-$4$ gerbe, matching the
  non-existence of a smooth $\ZZ/7$ CY$_3$ free quotient.
\end{itemize}

\subsection*{Cycle~5 --- consequence for the universal identity}

The universal Borcherds weight identity
$\kBKM(\Phi_N)=c_N(0)/2$ extends beyond the CHL slice $\{1,2,3,4,6\}$
to the full eight-form class $\{1,2,3,4,5,6,7,8\}$ \emph{iff} the
three non-CHL hosts $X_{5}^{\mathrm{BV}}$, $X_{7}^{\mathrm{?}}$,
$X_{8}^{\mathrm{MTW}}$ carry twisted DT zeta functions whose
$\Phi_N$-th metaplectic root matches the Borcherds lift of the
corresponding Gritsenko--Cl\'ery form. $N=8$ is the accessible
endpoint (weight $0$, Kummer-type host explicit); $N=5$ is the
Borcea--Voisin frontier (metaplectic square-root well-defined,
$(Z^X)^{-1/4}=\Phi_5$ a precise conjecture); $N=7$ remains open
pending identification of a singular CY$_3$ limit carrying the
spin-double-cover structure.

$\kcat$ on the non-CHL hosts is computed by Künneth on the
pre-quotient times the orbifold correction: $\kcat(X_8^{\mathrm{MTW}})
=\chi(\mathcal{O}_{\Kum^3(A)})/8+\kfib=1/8+\text{(McKay)}$, confirming
the non-Künneth-multiplicative behaviour on non-product CY$_3$.
$\kch$ is computed via $\Phi$ on the derived category of the smooth
resolution; the Hodge supertrace $\sum_q(-1)^q h^{0,q}$ agrees with
$c_N(0)/2$ modulo the $\kfib$ correction at the orbifold locus.

\subsection*{Primary sources}

Mongardi--Tari--Wandel, \emph{Symplectic involutions of K3$^{[n]}$
type and Kummer$^{n}$ type manifolds}, arXiv:1610.06216, 2018,
Thm~1.3; Gritsenko--Cl\'ery, \emph{The eight paramodular forms
with canonical diagonal divisor}, arXiv:0812.3962, 2008, Thm~1.2;
Nikulin, \emph{Finite groups acting on K3 surfaces}, Trudy Mosk.\
Mat.\ Obs.\ 38, 1980, Thm~4.5; Beauville, \emph{Variétés Kählériennes
dont la première classe de Chern est nulle}, J.~Diff.\ Geom.\ 18
(1983) 755; Borcea, \emph{K3 surfaces with involution and mirror
pairs of CY manifolds}, 1997; Voisin, \emph{Miroirs et involutions
sur les surfaces K3}, 1993; Borcherds, \emph{Automorphic forms
with singularities on Grassmannians}, Invent.\ Math.\ 132 (1998)
Thm~13.3; Oberdieck--Pixton arXiv:1808.03524 Conj.~A.
